\section{Preliminary User Study}\label{userstudy}
In this section, we investigate for a given topic which of the fake or real
headlines users are more interested in, and how often forward references are
found in interesting titles. Accordingly, we seek to test the following two hypotheses:

\noindent\textbf{H1}: \textit{Fake news headlines motivate user reading
	 interest more than real news headlines.}
  
\noindent\textbf{H2}: \textit{Forward references are commonly seen/used in headlines which interest users.}


We conducted the user study on both Chinese and English news to determine whether
forward references were used across languages.
For English headlines, we adopted FakeNewsNet~\cite{shu2018fakenewsnet}, which
contains fake and real news headlines about gossip and political news from
GossipCop and PolitiFact. Since the real and fake news in FakeNewsNet are not
paired up, we performed topical clustering to alleviate topical bias.  
For Chinese headlines, we directly leveraged news pairs labeled as
\textit{disagreed} in the WSDM fake news challenge
dataset,\footnote{\url{https://www.kaggle.com/c/fake-news-pair-classification-challenge}}
which contains one fake news headline and its corresponding verified news
headline. %In this way we avoid topical preferences.

We conducted the English user study using Amazon Mechanical Turk~\cite{amturk}.
Each pair was labeled by three turkers, whereas each Chinese pair was annotated by
five native speakers we recruited. 
To test H1, annotators chose which headline they wanted to read further, with
four options: \emph{first headline}, \emph{second headline}, \emph{both}, and
\emph{none}. News veracity was not revealed during the study. 
Results show that both Chinese and English readers prefer fake news headlines.
For Chinese headlines, 39.75\% of fake titles were judged to be more
interesting than the real ones, whereas only 23.60\% of real titles won. For
English headlines, the percentages are 34.57\% and 30.33\%, respectively. Note
that in English, we are comparing real news with fake news due to the scarcity
of paired verified and fake news data, whereas in Chinese, we are comparing
verified news with fake news. This could be why the preference for
real and fake news in English is closer than in Chinese. Even so, both Chinese and
English show with statistical significance (p-values far less than 0.05) that
readers prefer fake headlines. We report the complete distribution including
ties, as shown in Fig.~\ref{userstudy_dist}.
This result supports H1: fake news headlines motivate reading interest more
than real news headlines.
%: the tie ratio is 36.66\% and 35.10\% respectively.

\begin{figure}[h]
\centering
\subfigure[Chinese headlines]{\label{fig:a}\includegraphics[width=38mm]{figs/Chinese.png}}
\subfigure[English headlines]{\label{fig:b}\includegraphics[width=38mm]{figs/English.png}}
\caption{Reading preferences w.r.t.\ real news and fake news including ties.
The sample size is 8,424 / 6,497 for Chinese / English headlines.}
\label{userstudy_dist}
\end{figure}


To test H2, we randomly sampled 1,000 preferred and rejected headlines,
respectively, from the previous user study, and asked another set of three
annotators to label the FR type. 
%This question could have no answer, or more than one answer, as a title could be written without the use of FRs, or more than one FR type could be used to compose a title. 
%The definition of each FR type is listed in Appendix~\ref{app_B}. 
Results show that 73.48\% of Chinese and 85.32\% of English preferred headlines
utilizing FR techniques (at least one FR included in the headline), whereas in
rejected headlines, the ratio is 22.35\% / 17.72\%. This further supports
H2: FR is commonly used in interesting headlines. In conclusion, we found that
fake news headlines draw more reader interest, and the use
of FR techniques is a key part of what makes headlines intriguing.

% Below we provide the distribution of FR techniques being used in reader preferred headlines. Note that for Chinese headlines, we merged type~8 into
% type~1 as their definitions are similar in Chinese.\footnote{In Chinese,
% location adverbs are often regarded as a type of demonstrative pronoun.
% Also, as very few samples are of type~8, we treat it as a special case of type~1.} 
% Results show that the use of FR type depends on the language or culture.
% The distribution of forward-reference features in
% Chinese headlines is clearly more uneven. 
% This may be because the sources of news in our dataset are less diverse,
% resulting in more monotonic writing styles. 
% The majority of preferred titles utilize interrogatives (type~6) to lure readers
% to read the article to find the answer, while personal pronouns
% (type~2) appear less often.
% For English headlines, the types are observed roughly equally.
%In Chinese, utilizing interrogatives (type~6) to lure readers to read the article is most common, whereas in English all types are observed roughly equally.

% \begin{figure}[h]
% \centering
% \subfigure[Chinese headlines]{\label{fig:sub1}\includegraphics[scale=0.14]{figs/chpie.png}}
% \subfigure[English headlines]{\label{fig:sub2}\includegraphics[scale=0.15]{figs/enpie.png}}
% %\end{subfigure}
% \caption{Distribution of forward references in Chinese and English preferred headlines with 2,800/2,159 samples. Type~0 indicates headlines without 
% forward references.}
% \label{fig:test}
% \end{figure}


% \begin{figure}[h]
% \centering
% \subfigure[]{\label{fig:c}\includegraphics[width=38mm]{figs/FR_in_preferred.pdf}}
% \subfigure[]{\label{fig:d}\includegraphics[width=38mm]{figs/FR_in_dispreferred.pdf}}
% \caption{FR used in (a) preferred and (b) dispreferred headlines.}
% \label{FR_difference}
% \end{figure}

%When taking a closer look, we can easily find that the FR types in Chinese headlines are less evenly distributed. One possible reason is that the source of news in our dataset are less diverse, so the writing style is more monotonic

%It is also proven that FR is crucial in creating eye-catching titles to bait readers.

%In summary, our user study first discover the fake news headlines can attract users more than the verification headlines. Second, we verify FR is commonly seen or used in user-interested headlines, as high portion of preferred headlines contain at least one type of FR. This shows a direction for us to generate more user-desired headlines for verification news with forward-reference technique.


